% !TITLE 十一月四日风雨大作（其二）
% !AUTHOR 陆游
% !DYNASTY 两宋
% !GENRE 七言绝句
% !ORDER 17
% !EDITION 《陆游集》，中华书局（1976年版为候选）
% !EDITION_STATUS pending

\begin{poem}
僵卧孤村\textsuperscript{1}不自哀\textsuperscript{2}，尚思为国戍轮台\textsuperscript{3}。
夜阑\textsuperscript{4}卧听风吹雨，铁马冰河\textsuperscript{5}入梦来。
\end{poem}

\begin{annotations}
\notetext{1}{僵卧孤村：直挺挺地躺在孤寂的村庄里。僵卧，写年老体衰、卧病在床；孤村，写居处的荒僻。陆游写这首诗时六十八岁，已被罢官闲居在山阴（今浙江绍兴）老家。}
\notetext{2}{不自哀：不为自己哀伤。身体垮了，处境孤苦，他却不为这些难过。}
\notetext{3}{尚思为国戍轮台：还想着替国家去戍守轮台。尚，还；戍，防守；轮台是汉代西域的屯戍重地，在今新疆，这里借指北方的边防前线。}
\notetext{4}{夜阑：夜深，夜将尽的时候。}
\notetext{5}{铁马冰河：披着铁甲的战马和冰冻的河流，指北方严寒中的战场。风雨声在夜里听来像是战场上的动静，于是旧的军旅记忆就趁夜闯进了梦里。}
\end{annotations}

\begin{appreciation}
写这首诗时陆游六十八岁，已被罢官，闲居在山阴老家。诗题是风雨，诗里写的却是一具衰老的身体和一颗不肯衰老的心。

“僵卧孤村不自哀”——直挺挺地躺在荒僻的村子里，却不为自己难过半分。“僵卧”写身体，“孤村”写处境，“不自哀”写态度，七个字里三层意思一层比一层硬。

“尚思为国戍轮台”——还想着替国家去戍守轮台。轮台本是汉代设在西域的屯戍重地，在今天的新疆，这里借指北方的边防前线。全诗的力气都压在那个“尚”字上，身子已经僵卧在孤村而念头却还留在万里之外的轮台，两处隔得越远，这个“尚”字撑得越吃力。

“夜阑卧听风吹雨，铁马冰河入梦来”——夜将尽了，躺着听风吹雨打，披甲的战马与冰封的河流便一起闯进了梦里。前两句写的是醒着时的志，第三句才把风雨点出来，第四句一跌跌进梦里，由实入虚只在这半行之间。风雨声就是那座桥，雨点敲在屋瓦上，听得久了便是马蹄踏过冰河的声音。一个“入”字用得凶，梦是被这声音一把推进去的，容不得人慢慢做。

陆游一生写梦的诗极多，后人统计有近百首，梦里不是驰骋沙场就是收复中原。白天越是没有指望，夜里越是管不住自己的梦。同一个人，在《诉衷情》里说“心在天山，身老沧洲”，在这里说“铁马冰河入梦来”：前者把一生的话都说尽了，后者把话一直说到了梦里。

梁启超读陆游全集，下过一句断语：“亘古男儿一放翁。”这一夜他就这样听着满窗风雨睡去，铁马冰河还在梦里。醒来之后如何，诗没有说，窗外的风雨却还在下着。
\end{appreciation}
